% Standalone solution dossier for prop:a4-local-wellspecified.
\documentclass[../main.tex]{subfiles}
\begin{document}

% === SOLUTION HEADER (proof provenance) =====================================
% ledger-node: prop:a4-local-wellspecified
% refines: prop:a4-local-wellspecified
% bounded_by: obs:restricted-not-finite
% evidence_run: none
% checked_by: agent
% author: /root/a4_a5
% reviewer: /root/review_a4_a5
% dossier_editor: /root/review_a4_a5
% review: research/reviews/2026-08-21-a4-a5-independent-agent-audit.md
% date: 2026-08-21
% ============================================================================

\section*{Solution: well-specified local variational constants}

\begin{theorem}
Let $q_h\in\mathcal P_2(\R^d)$, $h\in\R^p$, be an identifiable family with $q_0=\pi$, so all
displayed means are finite. Assume that, for some $\rho_0>0$, the complete sublevel
$\{h:\KL(q_h\|\pi)\le\rho_0\}$ lies in a compact chart,
$h\mapsto\KL(q_h\|\pi)$ is continuous there, and $0$ is its unique zero. Suppose, uniformly
as $h\to0$,
\[
 2K(h)=h^\top Fh+O(\norm h^3),\quad
 D(h)=h^\top Gh+O(\norm h^3),\quad
 \E_{q_h}\theta-\E_\pi\theta=Bh+O(\norm h^2),
\]
where $K=\KL(q_h\|\pi)$, $D=W_2^2(q_h,\pi)$, and $F\succ0$. Then
\[
 C_{\mathcal Q,\rho}(\pi)
 =\lambda_{\max}(F^{-1/2}GF^{-1/2})+O(\sqrt\rho),
\]
and
\[
 C_{\mathrm{mean},\mathcal Q,\rho}(\pi)
 =\lambda_{\max}(F^{-1/2}B^\top BF^{-1/2})+O(\sqrt\rho).
\]
\end{theorem}

\begin{proof}
Continuity, compactness, and uniqueness imply that every sufficiently small KL sublevel lies in
an arbitrarily small neighbourhood of zero. In that neighbourhood, positive definiteness of
$F$ and the KL expansion give $K(h)\ge c\norm h^2$; hence $K(h)\le\rho$ implies
$\norm h=O(\sqrt\rho)$. On the punctured sublevel,
\[
 \frac{D(h)}{2K(h)}
 =\frac{h^\top Gh}{h^\top Fh}+O(\norm h)
 =\frac{h^\top Gh}{h^\top Fh}+O(\sqrt\rho)
\]
uniformly. Maximizing gives the upper bound. Arbitrarily small multiples of a top generalized
eigenvector belong to every positive sublevel and give the matching lower bound. Squaring the
mean expansion yields
$\norm{\E_{q_h}\theta-\E_\pi\theta}^2=h^\top B^\top Bh+O(\norm h^3)$, and the same argument
proves the mean statement.
\end{proof}

\paragraph{Obstruction respected.}
The compact, continuous, uniquely minimized sublevel assumption is explicit; a bare KL cutoff
over an unrestricted family is not used.

\end{document}
